% II. Анализ на алтернативите и обосновка на избора.
% ЗАДЪЛЖИТЕЛЕН за ОКС „Магистър“ (REQUIREMENTS.md А.3): „анализ на множество
% възможни решения и обосновка на направения конкретен избор“.
% Ориентировъчен обем: 6–8 стр.
%
% Съдържанието е изведено от port/cpp/PROFILE.md (§1, §4, §6, §8, §11) и
% port/cpp/docs/PUBLIC_API_DESIGN.md (§2.5, §4). Числата в §2.4 идват от
% преброяване в кодовата база — при повторно изброяване се обновяват.
\section{Анализ на възможните решения и обосновка на избора}
\label{sec:alternatives}

Всяко от разгледаните по-долу решения е било действителна възможност в момента на вземането
му, а не реторична алтернатива, въведена за да бъде отхвърлена. Разделът излага първо
критериите, после вариантите по петте определящи решения, и накрая --- обобщена оценка.
Където решението почива на измерване, а не на преценка, измерването е посочено.

\subsection{Критерии за оценка}
\label{subsec:criteria}

За да не бъде изборът въпрос на вкус, вариантите се съпоставят по пет предварително
изброени критерия:

\begin{enumerate}[topsep=2pt,itemsep=1pt,leftmargin=*,label=К\arabic*.]
\item \textbf{Проверимост} --- може ли резултатът от решението да бъде проверен
      автоматизирано, срещу външен източник на истина, без човешка преценка.
\item \textbf{Архитектурна вярност} --- запазва ли решението разделението, което прави
      описанието преносимо, или го размива.
\item \textbf{Независимост от среда за изпълнение} --- изисква ли решението
      допълнителна среда, вградена в крайния артефакт.
\item \textbf{Безопасност по конструкция} --- превръща ли решението цял клас грешки в
      невъзможни, вместо да разчита на дисциплина.
\item \textbf{Цена при добавяне на платформа} --- какво трябва да се напише отново при
      всеки нов платформен слой.
\end{enumerate}

Критерий К1 е поставен пръв съзнателно. Работата изследва твърдение за еквивалентност;
решение, което подобрява който и да е друг критерий за сметка на проверимостта, обезсмисля
изследването.

\subsection{Избор на целеви език и на езиковия стандарт}
\label{subsec:language}

\begin{table}[H]
\caption{Разгледани целеви езици по критериите К1--К5}
\label{tab:languages}
\centering\small
\begin{tabular}{@{}L{2.6cm}L{2.7cm}L{2.7cm}L{2.7cm}L{2.7cm}@{}}
\hline
\textbf{Критерий} & \textbf{C++26} & \textbf{Rust} & \textbf{Zig} & \textbf{C + обвивки} \\
\hline
К1 проверимост & пълна (пренасяне на тестовете е пряко) & пълна & пълна & пълна \\
К2 архитектурна вярност & висока & висока & висока & ниска (без обобщени типове) \\
К3 независимост от среда & да & да & да & да \\
К4 безопасност по конструкция & чрез типовата система & \emph{най-висока} & средна & ниска \\
К5 цена за платформа & \emph{най-ниска} (пряк достъп) & висока (обвързване) & висока & средна \\
\hline
\end{tabular}
\end{table}

Избран е \textbf{C++26} \tabref{tab:languages}. Rust печели по К4 --- проверката на
заемането прави безопасността свойство на езика, не на дисциплината --- но губи по К5,
който тук тежи повече: три от петте целеви инструментариума се обръщат от C++ без
междинен слой за маршалиране, а за Apple платформите смесеният език дава на една единица
за превод достъп едновременно до C++ и native обектния модел --- native обектите се държат
пряко в C++ клас, не се пренасят през граница. При Rust всеки от петте слоя би започнал с
изграждане на обвързване; а тъй като К4 се оказва постижим и в C++ чрез само преместваеми
типове (§\ref{subsec:ownership}), изборът пада върху C++26.

Изборът на \emph{конкретния стандарт} обаче не е формалност: задачата изисква (а)~работа
от изпълнение, преместена в превода; (б)~заместител на рефлексията, проверим от
компилатора; и (в)~модел на собственост, изразим в типовата система. Следващите
подпараграфи изброяват средствата, какво заместват и къде ефектът е реално измерим.

\subsubsection{Изчисления по време на компилация}

Оригиналът изгражда голяма част от структурите си при стартиране --- таблици за
преобразуване, идентификатори на типове, описатели на свойства --- работа, платена при
всяко пускане. В C++26 значителна част се измества в превода: \code{constexpr} и
\code{consteval} правят изчислението \emph{задължително} компилационно, не „по
възможност“ (\code{consteval} не компилира при извикване в контекст на изпълнение);
\code{if consteval} позволява различна реализация в двата режима без два отделни текста.
Резултатът е \term{константно инициализирани} константи и таблици, готови в артефакта, без
изграждане при пускане и без неопределения ред на инициализация между единиците за превод.

Делът на началното време, който отпада така, \emph{не} се твърди: измерването по H2b-1 не е
проведено (§\ref{subsec:cost}) --- механизмът е описан и предварително регистриран, но без
число, разликата между обосновка и реклама.

\subsubsection{Вграждане на описанието в превода}

Декларативното описание обичайно се чете от файл и се разбира при стартиране.
\code{\#embed} (C++26) прави съдържанието на файла част от единицата за превод, така че
разборът може да стане с \code{constexpr} код --- по време на компилация. Ползата е в
\emph{момента, в който грешката се открива}: неправилно описание престава да бъде отказ
при стартиране и става грешка при компилация.

Средството не е налично във всички поддържани компилатори; за една платформа това наложи
резервен път --- масив от байтове чрез отделна стъпка в изграждането
(§\ref{subsec:xaml_impl}), със същия резултат срещу цената на допълнителна зависимост.

\subsubsection{Концепции вместо проверки по време на изпълнение}

Концепциите (\code{concept}, \code{requires}) заместват два механизма на оригинала:
ограниченията върху обобщените типове, декларативни там, но проверявани отчасти при
изпълнение; и --- по-важно тук --- \term{незадължителните точки за разширение}. Базовият
клас на обработчиците открива при компилация дали производният предоставя дадена точка за
разширение чрез изискване за наличие на израз, така че производният изброява само
наистина предоставеното, а извикването на несъществуваща точка е грешка при компилация, не
безшумно нищо.

\subsubsection{Типизиран достъп вместо изтриване на типа}

Системата от свойства на оригинала стъпва върху общия базов тип на средата за изпълнение:
всяка стойност се \term{опакова}, а извличането проверява типа при изпълнение. Тук
описателят на свойство е шаблон по типа на стойността --- без опаковане, без информация за
типа при изпълнение, без заделяне за стойност, побираща се в самия обект.
Предимството между източниците на стойност е опаковано в едно 64-битово число, чийто
целочислен ред \emph{съвпада} с реда на предимството (§\ref{subsec:binding_design}):
разрешаването се свежда до едно целочислено сравнение.

\subsubsection{Наблюдения без копиране}

\code{std::span} и \code{std::string\_view} предават последователности и низове без
заделяне и копиране. В слой, в който имена на свойства се сравняват при всяко преминаване
през таблицата, разликата между сравнение на изглед и сравнение на собствен низ не е
козметична --- второто заделя памет на горещия път, изпълняван при всяка промяна.

\subsubsection{Детерминистично освобождаване}

Липсата на автоматично събиране на боклука е ограничение, но носи свойство, което
управляваната среда няма: \emph{моментът} на освобождаване се определя от обхвата, не от
събирач --- разрушителят освобождава native ресурса веднага.

Тук е нужна сдържаност: предварително регистрираната хипотеза \textbf{H2b-3} твърди липса
на разлика при установен режим на изобразяване (§\ref{subsec:preregistration}), защото при
малка купчина събирачът рядко надхвърля бюджета за кадър, а времето за кадър се изразходва
предимно в native инструментариума --- еднакъв и в двата случая. Очакваният измерим ефект е
в следата на паметта (H2a) и в началното време (H2b-1), не в честотата на кадрите.

\subsubsection{Само преместваеми типове}

Само преместваемите типове превръщат модела на собствеността от съглашение в проверимо от
компилатора свойство (§\ref{subsec:ownership}). Изтритият копиращ конструктор е старо
средство; новото е последователната поддръжка в стандартната библиотека, включително за
съхранение на извиквани обекти, което позволява един обработчик да улавя само преместваемо
състояние.

\subsubsection{Обобщение: къде езикът действително дава ефект}

\begin{center}
\begin{tabular}{@{}L{4.5cm}L{4.2cm}L{5.1cm}@{}}
\textbf{Средство} & \textbf{Какво замества} & \textbf{Очакван измерим ефект} \\
\hline
\code{constexpr}/\code{consteval} & изграждане при стартиране & начално време, размер \\
\code{\#embed} & четене и разбор при стартиране & начално време; момент на грешката \\
концепции & рефлексия, проверки при изпълнение & момент на грешката (не бързодействие) \\
типизиран описател & опаковане на стойности & следа на паметта, заделяния \\
опаковано предимство & сравнение на обекти & време при промяна на свойство \\
\code{span}/\code{string\_view} & копия на низове & заделяния по горещия път \\
детерминистично разрушаване & събиране на боклука & следа на паметта \\
само преместваеми типове & дисциплина по съглашение & коректност, не бързодействие \\
\end{tabular}
\end{center}

Третата колона е нарочно \emph{проверимо твърдение}, не реклама: три от осемте средства не
влияят на бързодействието --- плащат в коректност или в момента на откриване на грешка.
Смесването на двата вида полза е обичайният начин технически избор да бъде подкрепен с
числа, които не го подкрепят.

\textbf{Установено ограничение, непредвидено предварително.} Стандартната библиотека на
един компилатор не предоставя типа за само-преместваема функция, макар да е част от
стандарта; реализацията носи собствена, временна заместваща реализация --- показателно, че
изборът на стандарт не гарантира наличие на средствата му във всички компилатори, реална
цена на решението, не подробност.

\subsection{Стратегия за изграждане}
\label{subsec:port_strategy}

Разгледани са три стратегии:

\begin{enumerate}[topsep=2pt,itemsep=1pt,leftmargin=*]
\item \textbf{Автоматизирано преобразуване} на изходния текст на еталонната реализация.
      Отхвърлена: преобразуваният текст носи модела на паметта на изходния език (противоречи
      на К3 и К4), а резултатът не е проверим по К1 --- преобразуването не може да се твърди
      за коректно, освен чрез проверка, която връща задачата в началото.
\item \textbf{Ново прочитане} --- архитектура, вдъхновена от еталонната, но изградена
      наново. Отхвърлена по К1: без съответствие едно към едно между двете реализации няма
      общо описание, което и двете да изпълняват, тоест няма контролно условие за сравнение.
\item \textbf{Послойно изграждане, водено от тестове.} Избрана.
\end{enumerate}

Избраната стратегия работи слой по слой отдолу нагоре и за всяка компонента изисква
поведението да бъде \emph{прочетено}, не предположено, от три различни източника, всеки
отговарящ на различен въпрос:

\begin{center}
\begin{tabular}{@{}L{3.9cm}L{4.0cm}L{6.3cm}@{}}
\textbf{Въпрос} & \textbf{Източник} & \textbf{Какво дава} \\
\hline
Каква е формата? & описание на публичната повърхност & имена, сигнатури, наличност \\
Как трябва да се държи? & тестовете на еталона & стойности по подразбиране, гранични случаи, ред \\
Как е реализирано? & изходният текст на еталона & алгоритми, свързване с платформата \\
\end{tabular}
\end{center}

Разделянето не е педантизъм: формата не съдържа поведение, тестовете не съдържат алгоритъм,
а изходният текст не казва кое от направеното е договор и кое --- подробност. Замяната на
един източник с друг е обичайният начин да се получи правдоподобна, но невярна реализация.

Редът на разрешаване при противоречие е фиксиран предварително. Между \emph{форма} и
\emph{поведение} противоречие практически не възниква --- те описват различни неща.
Реалният случай е между \emph{изходния текст} на еталона и
\emph{действителното му поведение}: изходният текст е снимка от момент, а измерването е
срещу разпространената версия. Правилото, записано преди първото приложение, е:
\textbf{предимство има изобразеното} --- изходният текст остава източник за механизма, не
за стойност, която оттогава се е променила.

Първо приложение: изходният текст поставя долна граница от 44 единици за размера на
превключвател, а разпространената версия изобразява 18, без долна граница (измерено на две
платформи). Системата премахна границата, отбелязана като документирано отклонение с цитат
на двата източника; резултат --- два интерфейса преминаха от разминаване към точно
съвпадение.

\subsection{Модел на собствеността върху обектите}
\label{subsec:ownership}

Това е решението с най-големи последици --- засяга всеки ред от програмния интерфейс.
Дървото от изгледи съдържа обратни връзки родител\,$\leftrightarrow$\,дете и абонаменти за
събития, а средата на еталона ги разрешава чрез автоматично събиране на боклука, каквото тук
няма.

\subsubsection{Разгледани варианти}

\textbf{(а) Споделено притежаване навсякъде.} Най-близко до поведението на еталона и
най-лесно за писане. Отхвърлено, защото цикълът „контрола притежава обработчик, който
улавя контролата“ остава напълно безшумен: паметта никога не се освобождава, а
инструментите за откриване на течове мълчат, защото обектите са \emph{достижими} ---
просто взаимно.

\textbf{(б) Арена с индекси вместо указатели.} Решава цикъла по начина на изграждане и дава
компактно разположение в паметта. Отхвърлено по К2: заменя типовете с идентификатори, при
което програмният интерфейс губи типова информация и грешка в идентификатор става
невъзможна за откриване от компилатора.

\textbf{(в) Разпределение на ролите между четири средства.} Приетият вариант.

\subsubsection{Приетият модел}

Притежаващият манипулатор е \textbf{само преместваем}, наблюдаващият е копируем, а
съвместното притежаване се получава само чрез изричен метод. Следствието: най-честият
начин за цикъл --- улавяне на притежаващ манипулатор в обработчик, който самата контрола
притежава --- става \emph{грешка при компилация}, не безшумно и постоянно задържане на
паметта, което инструментите за откриване на течове не засичат.

\subsubsection{Емпирична обосновка}

Възражението срещу само преместваем манипулатор е ергономично: допълнително действие при
достъп до задържана контрола. Проверено е чрез преброяване в кодовата база, не преценено:

\begin{center}
\begin{tabular}{@{}L{9.2cm}r@{}}
\textbf{Наблюдение} & \textbf{Брой} \\
\hline
места за свързване към събитие, улавящи обхващащия обект & 322 от 392 \\
места, улавящи контрола чрез копие на притежаващ манипулатор & \textbf{0} \\
места с действителен достъп до задържана контрола, върху които пада цената & 111--130 \\
\end{tabular}
\end{center}

Предимството на копируемия манипулатор е предимство над случай, който на практика не се
среща, докато цената му --- безшумният цикъл --- се среща лесно. Преброяването обхваща
примерните приложения и тестовете в хранилището към момента на решението: измерва как
\emph{се пише} с интерфейса, не как \emph{би могло}.

\subsubsection{Съпоставка с наличната практика}

Съпоставката подкрепя избора и от друга посока: нито един зрял native инструментариум без
автоматично събиране на боклука не предоставя \emph{копируем притежаващ} манипулатор към
контрола --- всички въвеждат притежаване от родител към дете и отделен тип за наблюдаваща
препратка. Системите с копируем манипулатор са подпрени или от събирач, или от проверка на
заемането. Приносът тук не е в откриването на разпределението, а в това, че е
\emph{проверимо от компилатора}: същата препоръка, но като грешка при компилация вместо
бележка в документацията.

\subsection{Избор на платформени слоеве и на реда им}
\label{subsec:backends}

Реализирани са пет слоя. Редът, в който са изградени, е част от решението:

\begin{enumerate}[topsep=2pt,itemsep=1pt,leftmargin=*]
\item \textbf{Слоят без екран} --- първи. Няма native инструментариум и служи като
      определение за това какво изобщо трябва да реализира един платформен слой.
      Определящ по К1: докато не съществува, нищо от преносимата част не е автоматизирано
      проверимо.
\item \textbf{Настолният слой на разработващата машина} --- втори, защото доказва границата
      между преносимото и платформеното от край до край върху истински контроли.
\item \textbf{Мобилният и настолният вариант на същия инструментариум} --- трети,
      преизползва по-голямата част от втория.
\item \textbf{Настолният инструментариум на Windows} --- четвърти.
\item \textbf{Android} --- последен и най-скъп по К5: границата минава през чужд обектен
      модел с отделно управление на живота на препратките.
\end{enumerate}

Изборът кой слой да е втори следва от К5: смесеният език дава най-прекия достъп до native
обектите --- една единица за превод вижда двата свята, без междинно представяне.

\subsection{Подход към декларативното описание}
\label{subsec:xaml_choice}

Езикът не предоставя рефлексия, а еталонната реализация я използва точно за откриване на
обработчици, внедряване на зависимости и създаване на обекти от описанието. Изборът тук не
е стилов, а принудителен по същество:

\begin{enumerate}[topsep=2pt,itemsep=1pt,leftmargin=*]
\item \textbf{Тълкуване по време на изпълнение} --- изисква регистър на типовете, който
      така или иначе трябва да бъде попълнен явно; грешките в описанието се откриват при
      стартиране.
\item \textbf{Пораждане на код при компилация} --- грешките се откриват при компилация,
      няма разбор по време на изпълнение.
\item \textbf{Отказ от описание} --- само програмен интерфейс.
\end{enumerate}

Избран е вариант~(2) --- преобразуване по време на компилация --- с вариант~(1) като
резервен път там, където езиковото средство не е налично. Компромисът е конкретен: печели
се откриване на грешките при компилация и липса на разбор при стартиране; губи се
възможността описанието да се изменя без повторно компилиране --- съществено при инструмент
за бърза разработка, тук не е.

Явната регистрация е неизбежната цена на липсата на рефлексия: регистър на обработчиците
(кой договор с кой обект се обслужва), регистрация на свойствата (кое име към кой
описател) и регистър на преобразувателите (как текст става стойност от даден тип) --- всяко
в оригинала получено чрез обхождане на типовете при изпълнение.

Отбелязва се и капан от хода на работата: при самостоятелна регистрация (статичен обект,
вписващ се в регистъра при зареждане) свързващият редактор изхвърля единица за превод без
нито една препратка --- и регистраторът никога не се изпълнява. Решението: тези единици се
събират в обектна библиотека, която не подлежи на такова изхвърляне.

\subsection{Обобщение на направения избор}
\label{subsec:choice_summary}

\begin{table}[H]
\caption{Обобщение на петте решения по критериите К1--К5}
\label{tab:decisions}
\centering\small
\begin{tabular}{@{}L{3.4cm}L{4.6cm}L{6.0cm}@{}}
\hline
\textbf{Решение} & \textbf{Избран вариант} & \textbf{Решаващ критерий} \\
\hline
целеви език & C++26 & К5 --- пряк достъп до три от петте инструментариума \\
стратегия & послойно, водено от тестове & К1 --- само тя оставя общо описание, тоест контролно условие \\
собственост & четири роли, само преместваем притежател & К4 --- цикълът става грешка при компилация \\
платформени слоеве & пет, започвайки без екран & К1 --- без слоя без екран нищо не е проверимо автоматизирано \\
декларативен слой & преобразуване при компилация & К1 + К3 --- грешките се откриват по-рано, няма разбор при пускане \\
\hline
\end{tabular}
\end{table}

Две наблюдения върху \tabref{tab:decisions}. \textbf{Най-скъпото решение} е моделът на
собствеността: единственото, което засяга всеки ред от публичния интерфейс, и
единственото, чиято промяна по-късно би означавала пренаписване, не преработка.
\textbf{Най-определящото за останалите} е слоят без екран --- докато не съществува, никое
от другите решения не може да бъде проверено, а решение, което не може да бъде проверено,
на практика не е взето.

Забележимо е, че К1 е решаващият критерий в три от петте реда --- не съвпадение: работата
изследва твърдение за еквивалентност, и решение, което затруднява проверката, отнема от нея
повече, отколкото дава.
